\documentclass[11pt]{article}
\usepackage[margin=1in]{geometry}
\usepackage{amsmath,amssymb,amsthm,mathtools}
\usepackage[hidelinks]{hyperref}
\newtheorem{lemma}{Lemma}
\newtheorem{proposition}{Proposition}
\theoremstyle{remark}
\newtheorem{remark}{Remark}
\newcommand{\Tr}{\operatorname{Tr}}
\newcommand{\HD}{\mathcal H_D}
\newcommand{\one}{\mathbb 1}

\title{Red-Team Comment on Chandrasekaran, arXiv:2611.34924v1:\\The Finite-Dimensionality Argument of \S5.2 Rests on an Unproved Gluing/Normalizability Premise}
\author{Justin\\\small Nous Enterprises LLC}
\date{27 November 2026 --- red-team technical comment, v1.4}

\begin{document}
\maketitle

\section*{Thesis}
Given the operator bound (5.30c), the \S5.2 conclusion $\dim\HD<\infty$ follows in one
line from a finite uncut norm together with norm-preserving gluing, applied to the
identity insertion. The later boost--cycle sign argument (5.31--5.38) is therefore
logically redundant once that premise is granted. The manuscript justifies the gluing
step by diffeomorphism invariance and Eq.~(4.29). Eq.~(4.29) controls open matrix
elements between fixed normalizable vectors. It does not control the closure of the
doubled preparation. Diffeomorphism invariance under changes of cut is condition~1 of
the paper's axiomatic Appendix~E. Its realization (E.15a),
$\zeta_L(\mathrm{cl}_\Sigma a)=\Tr_{H_\Sigma}\hat a$, \emph{asserts} that closing a
preparation equals the trace over states on the cut, rather than deriving it. Applied to
$\hat a=T$, this makes the closure equal to $\Tr T$. Finiteness of such closures is
condition~5, and the manuscript verifies condition~5 for $\HD$ in (E.17) by invoking
$d_D<\infty$ from \S5.2 (via (5.39c)). Hence:

\begin{quote}
\emph{Within the manuscript, the finiteness that enters the \S5.2 gluing step is either
postulated (Appendix~E, conditions~1 and 5 with (E.15a)) or established using \S5.2
itself (E.17). No derivation from the reduced dynamics that is independent of \S5.2 was
located. Once that premise is granted, finite dimension follows without the
boost--cycle argument.}
\end{quote}

The paper's own JT prologue sharpens this. There, under the same diffeomorphism-invariant
gluing, closing the preparation gives $\Tr W=\infty$ for an operator bounded below on an
infinite orthonormal family (2.17b). So diffeomorphism invariance does not by itself
supply a finite uncut norm, even within the manuscript (Section~\ref{sec:more}).

This is a missing proof obligation. It is not a refutation of the proof, and it is not
a claim that $\HD$ is infinite dimensional. Section~\ref{sec:more} also shows that the
headline log-coefficient match with Sen is an algebraic identity of the one-loop
heat-kernel input, and so carries no information about the null-cone quantization. It
also notes limits (F5--F7) on how far the paper's other results support its broader
claims. The graviton-pair amplitude is a sound order-of-limits example, but it is a
free-field series and does not by itself bear on interacting UV behavior. The
semiclassical limit is tested through one loop, which does not probe the order where
pure gravity's known nonrenormalizability appears. The type~I structure follows from the
choice of a pure state, and the paper leaves open whether that choice matters.

\section{Premises extracted from the manuscript}
\begin{itemize}
  \item[(S1)] (5.23), (5.25), (5.30c): $T\coloneqq C_b^\dagger C_b=e^{-\hat s}f_b(\hat s)$
  on $\HD$, with $c\,\one\le T\le\one$ and $c=e^{-s_*}m_{b,D}>0$. The positivity of
  $c$ uses the focusing-based area bound (5.21), (5.30a), so that upstream focusing input
  is \emph{not} redundant.
  \item[(S2)] Doubled cut preparation. For a finite-rank basis projection $P$,
  $\Phi_P=\sum_{i\in P}C_b e_i\otimes\bar e_i$, with $\|\Phi_P\|^2=\Tr(PT)$.
  \item[(S3)] $Z_P[X]=\Tr(PT)\,\langle X\rangle_P$ (5.38).
  \item[(S4)] Gluing (pp.~35, 42--43): as $P\to\one$ through complete-cycle unions,
  $Z_P[X]\to Z[X]=\|C_b\|^2_{\rm uncut}\langle X\rangle_{C_b}$. Justification offered:
  (4.29) and diffeomorphism invariance (\S3.6).
  \item[(S5)] $X=\Delta_\eta\hat O-\tfrac{\delta_\eta}{2}\one$, with
  $\langle X\rangle_P\le0$ on every cut and $\langle X\rangle_{C_b}>0$ (5.37).
\end{itemize}

\section{Results}
\begin{lemma}[Trace--dimension equivalence]\label{lem:equiv}
Under (S1), every finite orthonormal family $F$ satisfies $c|F|\le\Tr(P_FT)\le|F|$.
With $\Tr T\coloneqq\sup_F\Tr(P_FT)\in[0,\infty]$, we get $\Tr T<\infty$ if and only if
$\dim\HD<\infty$, and $c\dim\HD\le\Tr T\le\dim\HD$. No separability is assumed.
\end{lemma}

\begin{proposition}[Identity insertion]\label{prop:one}
If (S4) holds for $X=\one$ with $N\coloneqq\|C_b\|^2_{\rm uncut}<\infty$, then
$\dim\HD\le N/c$. The boost--cycle argument (5.31--5.38) is not needed.
\end{proposition}
\begin{proof}
$Z_P[\one]=\Tr(PT)$ increases monotonically to $N$, and Lemma~\ref{lem:equiv} then
bounds $|F|\le N/c$.
\end{proof}

\begin{remark}[Scope]
Only \emph{finite trace} is equivalent to finite dimension. The full content of (S4)
is stronger. It identifies the cut limit with a specific gravitational amplitude, with
a prescribed measure, normalization, and insertions. Equivalence between a premise and
the conclusion is not, by itself, circularity. An independent proof of finite uncut
norm plus a valid gluing theorem would establish the result through
Prop.~\ref{prop:one}.
\end{remark}

\begin{proposition}[What fails without the premise]\label{prop:none}
Suppose $\dim\HD=\infty$. Then:
\begin{enumerate}
  \item $\|\Phi_P\|^2\ge c\,\mathrm{rank}\,P\to\infty$. The \emph{cut} preparations
  diverge, so (4.29), whose bound is $\le\|\phi\|\|\psi\|$ for fixed vectors, does not
  apply to their closure. This says nothing about whether an \emph{uncut} object is
  finite, fails to exist, or is simply not the limit of the cuts. That identification is
  exactly what is disputed.
  \item For the signed insertion, $Z_P[X]=Z_P[\Delta_\eta\hat O]-\tfrac{\delta_\eta}{2}
  \Tr(PT)$ is a difference of two divergent quantities. No trace or cancellation
  prescription is supplied for the limit.
  \item The normalized one-sided densities $\rho_P=PTP/\Tr(PT)$ need not determine the
  two-sided expectation $\langle X\rangle$. Even in $\mathcal B(\ell^2)$, normalized
  finite-rank identity densities have only singular weak-$*$ cluster points. The failure
  therefore occurs within type~I and does not by itself diagnose a non-type-I algebra.
\end{enumerate}
\end{proposition}

\paragraph{Separation model (due to the Codex review; verified by hand).}
Let $H=\ell^2(\mathbb N)\oplus\mathbb C f$ with $\hat s e_n=0$ and $\hat sf=f$. Take
$T=q\one\oplus qe^{-1}$, $q=e^{-1/8}$, $C=T^{1/2}$, $u=e_1\otimes\bar e_1$,
$v=f\otimes\bar f$, $O=-|u\rangle\langle v|-|v\rangle\langle u|$, and
$\chi=(u+v)/\sqrt2$. A one-sided boost multiplies $v$ by $e^{i\theta}$, which gives
$\langle\Delta O\rangle_\chi=\delta\coloneqq1-\cos\theta$. Put $\{f,e_1,e_2,e_3\}$ in
one 4-cycle, with $k$ further 4-cycles. Then
\[
\langle\Delta O\rangle_P=\frac{2e^{-1/2}}{3+e^{-1}+4k}\,\delta<\frac\delta2,
\]
so every cut has negative unnormalized amplitude while $\langle X\rangle_\chi>0$. This is
consistent because $\chi$ is not the cut limit of $C$. The model is not gravitational.
It isolates the missing identification step.

\section{Where the gluing step could be grounded}
\begin{center}
\begin{tabular}{p{0.27\textwidth}p{0.63\textwidth}}
\hline
Candidate & Status\\\hline
Eq.~(4.29)--(4.31) & Controls open matrix elements and finite subdivisions with fixed external vectors. Insufficient for trace closure.\\
\S3.6 diffeo.\ invariance & Asserted. Not a quantitative convergence theorem.\\
App.~E, cond.~1, (E.15a) & Closure $=$ trace on the cut is asserted (``we can evaluate the same amplitude by closing \ldots\ or by summing over the states''). Applied to $T$, it gives closure $=\Tr T$.\\
App.~E, cond.~5, (E.4--E.5) & Finite closures are required axiomatically and verified in (E.17) using $d_D<\infty$ from \S5.2 via (5.39c). Citing this would close the cycle.\\
Introduction (p.~5) & States that the Lorentzian axioms match the Euclidean ones ``when the Euclidean cylinder has finite trace.'' Finite trace appears here as a condition, not a result.\\
\S5.4, (5.71--5.74), p.~50 & Factorization and Hilbert--Schmidt structure of connected fillings, developed \emph{after} \S5.2. Not an independent input.\\
\hline
\end{tabular}
\end{center}

\section{Further obligations in the sign argument (from the Codex review)}
These are independent of the trace observation. They would matter even if the gluing
step were repaired.
\begin{enumerate}
  \item \textbf{Approximation control.} The vacuum response (5.28--5.29, D.15--D.20)
  and (D.36) are computed in the linearized shear description. The strict exact
  inequality (5.37b) needs a bound on the full remainder that is smaller than
  $\delta_\eta/2$ in a single compatible parameter regime. Power counting in
  $\ell_P^2/(\Delta u)^2$ does not supply this.
  \item \textbf{State identification.} The filter $e^{-\hat s}f_b(\hat s)$ distinguishes
  areas, not radiative states at fixed area. Assigning the smooth-background excitation
  bound (5.36b) to $C_b$ therefore needs a selection or dominance argument.
  \item \textbf{Exact self-adjointness.} $\sin(j_L\hat\sigma_{f_L})\sin(j_R\hat\sigma_{f_R})$
  is self-adjoint only if the factors commute. Commutativity is shown at linearized
  order. Gravitationally dressed observables at spacelike separation need not commute,
  so this is not merely cosmetic. Symmetrizing the product is a candidate repair.
  \item \textbf{Parameter hierarchy.} $\ell_P\ll\Delta u\ll d\ll B\ll b\ll R_c$ must fit
  inside a fixed diamond whose area is bounded.
  \item \textbf{Acknowledged heuristic upstream.} Footnote~18 (after (5.16)) calls the
  transfer of a classical tail estimate to the exact matrix element a heuristic. This
  bears on the projected preparation used in (5.20).
\end{enumerate}

\section{Further findings}\label{sec:more}

\paragraph{F1. The JT prologue is an in-paper test of the gluing premise.}
In \S2 (pp.~10--11) the paper glues the JT preparation $C_D$ to its time reflection and
obtains $\langle\Psi_m|W|\Psi_n\rangle=e^{-2\pi(K_L+K_R)/\hbar}\delta_{mn}$ for an
infinite orthonormal family. Hence $\Tr_{\HD}W\ge N e^{-2\pi(K_L+K_R)/\hbar}$ for
every $N$, and ``$\Tr W=\infty$ and $W/\Tr W$ cannot define a density matrix''
(2.17a--b). This is exactly the structure of (S1), a positive lower bound on an infinite
family, inside a framework with the same diffeomorphism-invariant gluing and the same
appeal to summing over states on the cut. The justification offered for the 4D gluing
step (diffeomorphism invariance; (4.29)) is not specific to four dimensions. If it
implied a finite uncut norm, it would imply one in JT, contradicting (2.17b).

The paper attributes the 4D difference to focusing, which restricts radiative states.
In \S5.2, however, focusing enters only through the sign comparison (5.37--5.38). That
comparison needs $Z[X]$ to be a definite finite quantity. With $Z[\one]=\infty$, as in
JT, $Z[X]=Z[\Delta_\eta\hat O]-\tfrac{\delta_\eta}{2}Z[\one]$ has no prescribed value. The
abstract itself credits the result to ``diffeomorphism invariance of the path
integral.'' The JT case shows that this is not sufficient. A finite 4D uncut norm is an
additional input that is not established before \S5.2.

\paragraph{F2. The $\tfrac{106}{45}$ match is an algebraic identity of the input.}
Let the one-loop computation give $\Tr(C_b^\dagger C_b)\propto s_b^{\,p}$ (5.67c). Then
(5.66) gives $Z(\ell)\propto s_*^{\,p+1/2}\propto\ell^{\,2p+1}$, and (5.58) gives
$q/(1-q)=2p+1$, i.e.\ $q=(p+\tfrac12)/(p+1)$. The fixed-mass coefficient (5.69) is then
\[
\frac{q}{2(1-q)}-\frac12 \;=\; p\qquad\text{identically in }p .
\]
(Verified symbolically, along with $q(106/45)=257/302$. The normalization in (5.66),
\[\int_0^\infty e^{-(s-t)^2/4t}\,(4\pi t)^{-1/2}\,dt=1,\]
is exact, so the paper is correct there.) The clock and ensemble machinery therefore maps the heat-kernel exponent back to
itself, and the reported coefficient \emph{is} the one-loop conical-defect coefficient
$p=\tfrac{106}{45}$. Sen's mixed-ensemble result for pure 4D gravity is
$C_{\rm local}\ln a=\tfrac{212}{45}\ln a=\tfrac{106}{45}\ln A$. In this case it also
reduces to the local Seeley--DeWitt term \cite[Eqs.~(1.2), (1.4)]{Sen}. The
agreement is therefore guaranteed by both calculations using the same local heat-kernel
coefficient. It tests neither finite dimensionality nor the null-cone construction,
and in four dimensions it would hold for any matter content whose heat-kernel coefficient is used as input. The ``Lorentzian'' route also uses an
opening angle $4\pi$ at $S$ (Fig.~33) and an analytic continuation of the boost to a
$2\pi$ rotation, which is the conical-singularity method underlying the Euclidean
calculation. It is not an independent method.

\paragraph{F3. The Bekenstein--Hawking term is fixed by the corner weight.}
$S(\rho_D)=\langle\hat A\rangle/4G_N\hbar+\log Z_{D,1}$ (5.50) is the Gibbs identity for
$\rho_D\propto e^{-\hat s}$. The weight $e^{-s_i}$ comes from continuing the corner boost
through two half turns (\S5.3, Fig.~32), and the paper says plainly that this
``boost weight yields the Bekenstein Hawking term.'' The state count (5.64),
$\dim\sim e^{as_*}$, then follows from that weight together with $\log Z=o(s_*)$ and
finite dimensionality. This is internally consistent. It is not, however, an
independent microscopic derivation of $A/4G$, and it inherits the \S5.2 conditional.

\paragraph{F4. A second instance of the premise at the AdS boundary.}
The finiteness of $Z_\beta=\Tr(C^\dagger C)$ (5.44--5.45) is derived by requiring that
closing the joint cut ``must yield zero.'' The argument uses $d_D<\infty$ from \S5.2,
and it again presumes that the closed amplitude is the limit of partial sums.
Discreteness of the AdS spectrum (5.46), the type~I boundary algebra (5.47--5.48), and
the Hartle--Hawking state (5.49) inherit both dependencies.

\paragraph{F5. The \S4.4 example illustrates an order-of-limits effect, but does not establish UV finiteness of interacting amplitudes.}
The paper presents this calculation as an illustration of exchanging limits, and as such
it is correct. The concern is only with reading it as evidence for broader UV claims.
The ``nonperturbative'' amplitude (4.33) is
$\sum_k 2(\Delta q_k)^4e^{-2ik/s}=\mathrm{Li}_2(e^{-2i/s})/(2s^2)$, with
$(\Delta q_k)^2=1/(2sk)$ (verified numerically). Three features limit what it shows.
(i) Absolute convergence comes from the $k^{-2}$ Gaussian weight, not from gravity. The
sum converges for every $s$ even without the phase. At fixed area the amplitude itself
vanishes as $G_N\to0$, while the rescaled quantity $s^2A(s)\to\zeta(2)/2$. (ii) The phase $e^{-ik/s}$ per graviton is \emph{chosen}: ``we choose
this relative phase \ldots\ by rotating the modes with a boundary unitary,'' an
``affine shift'' motivated by analogy with a string worldsheet (p.~34). It is not derived
from dynamics. (iii) The $\log N$ in (4.34) is the non-analyticity of a Fourier series
at the edge of its domain:
$\mathrm{Li}_2(e^{-it})-\pi^2/6=-it\,(1-\log(it))+O(t^2)$. Expanding term by term in
$t=2/s\propto G_N$ produces $\log N$, and resumming replaces it with $\log s$. This
reflects a mode-number crossover at $k\sim s$ set by the chosen phase. The paper does not
show that it corresponds to a physical (e.g.\ Planck-momentum) cutoff. The computation involves a
Gaussian state and quadratic insertions, so it has no interaction vertices and does not
engage the two-loop $R^3$ divergence~\cite{GS} that \S4.4 cites as its
motivation. The asymptotics
$A(s)=\zeta(2)/(2s^2)+i(\log(2i/s)-1)/s^3+O(s^{-4})$ show that $A$ is finite for
every $G_N>0$ and continuous at $G_N=0$, and that the divergent third derivative in (4.34)
is consistent with this. Eq.~(4.30) is Cauchy--Schwarz, \emph{given} a Hilbert-space resolution of the identity,
normalizable fixed external states, and unitary evolution. Once those are in place, the
bound is generic, so absolute convergence of intermediate sums is not specific evidence
of UV completeness.

\paragraph{F6. One-loop semiclassical agreement does not establish an all-orders UV completion.}
\S6.1 obtains the Einstein equations from stationarity of the reduced GR action in the
$\hbar\to0$ limit of coherent states, which is the correspondence principle for a
quantization of GR's own phase space. \S6.2 and App.~D.1 show agreement with constrained
GR's \emph{one-loop} generating functional, at fixed smearing ((6.14): ``with smearing
held fixed''). In the familiar setting, pure gravity with $\Lambda=0$ has an on-shell one-loop
S-matrix free of divergences, up to topological terms~\cite{tHV}. The non-renormalizable
$R^3$ counterterm first appears at two loops~\cite{GS}. With $\Lambda\neq0$, one-loop
on-shell divergences proportional to the Euler density and $\Lambda^2$ survive (see,
e.g.,~\cite{BPS}). That on-shell result is
narrower than what the paper examines: \S6 allows $\Lambda\neq0$, and it works with a
generating functional and coincident-point quantities at fixed smearing. The one-loop match is a substantive low-energy consistency check. It does not show UV
finiteness to all orders, and it does not compute a two-loop amplitude. The natural next
test is to show how the framework reproduces or resolves the known two-loop low-energy
behavior while defining its full amplitude. Its perturbative $R^3$ coefficient need not
vanish, since a finite nonperturbative amplitude can have a singular expansion, which is
the mechanism of (4.33)--(4.34). We also ask the author to state the regulator,
renormalization, and boundary conditions under which (D.2), (D.12), and (6.14) exist, and
whether the comparison is made before or after renormalization.

\paragraph{F7. The type~I bicommutant follows from a pure-state extension; representation independence remains open.}
In \S4.2 the paper extends the incoming Gaussian state to a \emph{pure} state $\omega$
on the full algebra, citing~\cite{KS}, and concludes
$\pi_\omega(\mathcal A)''=\mathcal B(\mathcal H_\omega)$ (4.17). For any pure state the
GNS representation is irreducible, so its commutant is scalar and its bicommutant (weak
closure) is $\mathcal B(\mathcal H_\omega)$. The norm-closed image $\pi_\omega(\mathcal A)$
may still be a proper subalgebra. This holds for every C$^*$-algebra. When the algebra is
a type~III von Neumann algebra considered only as a C$^*$-algebra, such a representation
is generally non-normal with respect to the original one and does not change that
algebra's type. The type~I bicommutant (4.17) therefore follows from choosing a pure-state
GNS representation. It is not, by itself, a consequence of gravitational dynamics. Pure
extensions need not be unique. For example, with $\mathcal A=M_2(\mathbb C)\oplus
M_3(\mathbb C)$ and the scalar subalgebra, the unique state on the subalgebra has pure
extensions on either summand. Both have type~I bicommutants, but they are inequivalent,
of dimensions 2 and 3. This toy does not model the incoming algebra. It shows only that
an extension theorem cannot by itself supply representation independence. Footnote~20
defers this question. The relevant condition is not uniqueness of $\omega$ but unitary
equivalence that intertwines the clock, area, and preparation operators. Whether $d_D$, the area spectrum, or the entropy in \S5.3
depend on $\omega$ is therefore an open possibility. It has not been demonstrated in
either direction. The manuscript cites the author's earlier work on gravitational algebras of
type~II (refs.~[18, 19] of~\cite{VC}), but it does not explain how a type~I factor obtained by choosing a pure
state relates to those type~II results. The constructions differ, so the coexistence of
the two is not a contradiction. The point is that an explanation would be useful. The JT prologue, where the paper finds a
type~I$_\infty$ factor with $\Tr W=\infty$ (F1), shows that this type~I structure is
compatible with infinite closure. That alone does not show that the 4D representation
choice causes (5.38) to fail.

\paragraph{Checked and consistent.} Footnote~23's $q=2x^2/(1+x^2)$ agrees with the
one-loop $q=257/302$. The ratio $\ell(D)/r_h$ is not free; it is fixed by the ensemble
($x^2=\langle s\rangle/s_*=257/347$).

\section{Machine-checked component}
\texttt{TraceDim.lean} compiled unchanged under Lean~4.35.0, Mathlib
\texttt{56a0058}. All three declarations depend only on \texttt{propext},
\texttt{Classical.choice}, and \texttt{Quot.sound}, with no \texttt{sorryAx}. SHA-256
\texttt{789ba2c8\ldots664893} matches between authoring and verification (Codex
package). Scope: linear growth of finite orthonormal partial sums under a quadratic-form
lower bound, and finiteness of the index type of an orthonormal family under the
\emph{assumed} bound \texttt{hN}. Not formalized: the conclusion
\texttt{FiniteDimensional} (this needs a Hilbert basis), the converse direction, the
upper bound $T\le\one$, and any physical premise.

\section{Assessment}
The elementary operator deductions hold. The manuscript's $\dim\HD<\infty$ is
conditional on a gluing/normalizability theorem that is not supplied independently.
The sign argument carries further obligations of approximation control,
state identification, and exact self-adjointness. Downstream results (discrete area
spectrum, (5.39a), and the \S5.3 counting) inherit the conditional. The
$\tfrac{106}{45}$ match is an identity of the one-loop heat-kernel input (F2) and is no
evidence for finiteness.

\paragraph{Cleanest repair.} Prove $C_b\in\mathcal S_2$ (Hilbert--Schmidt), or a finite
norm-preserving closure for this specific preparation, from the reduced dynamics and
without using \S5.2 or Appendix~E. Prop.~\ref{prop:one} then completes the argument.
Alternatively, state finite closure as an explicit axiom and demonstrate a nontrivial
realization of it.

\paragraph{Scope.} This is a rapid red-team comment, not a full review.
Appendix~D.3, the reduced quantization of \S\S3--4, and the entropy calculation beyond
the coefficient check were not audited, and \S6 was examined only for scope (F6). Nothing here shows that $\HD$ is
infinite dimensional or that the framework is inconsistent. Corrections from the author
are invited and will be incorporated in a revised version.

\paragraph{Contributions and AI use.} The author directed the review, selected its
targets, and takes responsibility for all claims. AI systems were used as research
tools, not authors. ChatGPT (OpenAI) produced an initial triage that flagged \S5.2.
Claude Opus~5.5 (Anthropic) reduced the objection to the identity-insertion argument,
identified F1--F7,
drafted the note, and wrote \texttt{TraceDim.lean}. Claude's analysis was primed by the
ChatGPT triage. OpenAI Codex then reviewed the draft adversarially. It corrected several
of its claims, identified the (E.17) dependency and the obligations in Section~4,
constructed the separation model, and compiled the Lean file. A second Codex pass on
v1.2 narrowed the wording of F5--F7, supplied the $M_2\oplus M_3$ illustration, and
contributed the asserting tail-bound check for (4.33). Because each system saw
earlier outputs, their agreement is not independent replication. Every manuscript
citation in this note was checked against the v1 PDF.

\paragraph{Artifacts.} \texttt{TraceDim.lean}, the verification log, the
finite-matrix checks, and the separation-model script are included as supplementary
files.

\begin{thebibliography}{9}
\bibitem{VC} V.~Chandrasekaran, \emph{An algebraic quantization of gravity from the null cone}, arXiv:2611.34924v1 [hep-th] (2026).
\bibitem{tHV} G.~'t~Hooft and M.~Veltman, \emph{One-loop divergencies in the theory of gravitation}, Ann.\ Inst.\ H.\ Poincar\'e A \textbf{20} (1974) 69--94.
\bibitem{GS} M.H.~Goroff and A.~Sagnotti, \emph{The ultraviolet behavior of Einstein gravity}, Nucl.\ Phys.\ B \textbf{266} (1986) 709.
\bibitem{BPS} G.~de~Berredo-Peixoto, A.B.~Penna-Firme, and I.L.~Shapiro, \emph{1-loop divergences of quantum gravity using conformal parametrization}, doi:10.1142/S0217732300002929, arXiv:gr-qc/0103043.
\bibitem{Sen} A.~Sen, \emph{Logarithmic corrections to Schwarzschild and other non-extremal black hole entropy in different dimensions}, arXiv:1205.0971.
\bibitem{KS} R.V.~Kadison and I.M.~Singer, \emph{Extensions of pure states}, Amer.\ J.\ Math.\ \textbf{81} (1959) 383--400.
\end{thebibliography}

\end{document}
